\chapter{Unit 5: Testing, Validation and Customer Experience}
\label{chap:unit5_qb}

\section*{Practice Question Set (30 MCQs)}
\addcontentsline{toc}{section}{Practice Question Set (30 MCQs)}

\mcqitem{1}{What is the primary purpose of user testing in Design Thinking?}{To prove to management that the design is flawless}{To observe real user behavior and identify usability defects, confusion, and invalid assumptions}{To grade the intelligence of the user}{To sell subscriptions to test participants}

\mcqitem{2}{What is the fundamental difference between Usability Testing and Design Validation?}{Usability testing checks if the interface works efficiently; validation checks if it solves the right user problem}{Usability testing is done by developers; validation is done by accountants}{Usability testing requires paper; validation requires code}{There is no difference; they are exact synonyms}

\mcqitem{3}{Which of the following best exemplifies the "Show, Don't Tell" testing approach?}{Presenting a 30-minute PowerPoint pitch explaining why the app is easy to use}{Handing the user a prototype with a task scenario and observing their actions silently}{Writing an instruction manual and reading it aloud to the user}{Showing a promotional video of the product}

\mcqitem{4}{During a neutral user testing session, what should the facilitator say when a user asks, "Should I click this button?"}{"Yes, that is the correct button!"}{"What do you think will happen if you click it?"}{"No, don't click that, it's wrong."}{"Let me click it for you."}

\mcqitem{5}{Why is it important to tell the participant: "We are testing the product, not you"?}{To protect the company from legal liability}{To reduce participant anxiety and ensure they do not feel at fault when design errors occur}{To prevent the participant from asking for money}{Because participants are not allowed to make mistakes}

\mcqitem{6}{What makes a user testing Task Scenario effective?}{It specifies every button and menu name in exact order}{It provides a realistic goal and context without dictating specific interface steps}{It has no defined end goal}{It can be completed in less than one second}

\mcqitem{7}{Which of the following is a POORLY constructed task scenario?}{"Find a pair of running shoes under Rs.~2,000 and proceed to the shipping screen."}{"Click the search bar, type 'shoes', click filter, select 'under 2000', and tap 'Add to Cart'."}{"You need to book a flight to Mumbai for next Monday morning. Show how you would do that."}{"Check your account balance and download last month's statement."}

\mcqitem{8}{What is the "Think-Aloud Protocol" in usability testing?}{The facilitator constantly explains what the software is doing}{The participant verbalizes their thoughts, expectations, and hesitations aloud while performing the task}{All observers in the room shout feedback at the participant}{The participant recites the user manual from memory}

\mcqitem{9}{Which quadrant of the Feedback Capture Matrix records constructive criticisms and unmet user wishes?}{Likes (+)}{Wishes ($\Delta$)}{Questions (?)}{Ideas (!)}

\mcqitem{10}{In a Feedback Capture Matrix, where does a user's statement *"Is international shipping included in this price?"* belong?}{Likes (+)}{Wishes ($\Delta$)}{Questions (?)}{Ideas (!)}

\mcqitem{11}{How does Customer Experience (CX) differ from Usability (UI/UX)?}{Usability is broader than CX}{Usability evaluates specific interface interactions; CX encompasses the holistic end-to-end relationship across all touchpoints}{Usability only applies to websites; CX only applies to physical retail stores}{Usability is measured in revenue; CX is measured in clicks}

\mcqitem{12}{Which of the following is an example of a HUMAN touchpoint in a service journey?}{An automated billing email}{A delivery courier handing a package to the customer}{The website homepage}{The cardboard product packaging}

\mcqitem{13}{In Customer Journey Mapping, what are the primary layers mapped across stages?}{User Actions, Thoughts/Emotions, Touchpoints, Pain Points, and Opportunities}{HTML code, CSS stylesheets, JavaScript files, and Server logs}{Profit margins, Depreciation rates, and Tax brackets}{Employee salaries and Office desk locations}

\mcqitem{14}{What does the "P" stand for in the WCAG accessibility POUR principles?}{Predictable}{Perceivable}{Portable}{Practical}

\mcqitem{15}{According to WCAG 2.1 guidelines, what is the minimum color contrast ratio required for standard body text (Level AA)?}{1.5:1}{3.0:1}{4.5:1}{10:1}

\mcqitem{16}{What accessibility feature ensures that visually impaired users using screen readers can understand images?}{High-resolution 4K images}{Descriptive Alternative Text (`alt` attribute)}{Flashing neon borders}{Background audio music}

\mcqitem{17}{Which of the following violates the "Operable" principle of web accessibility?}{An interface that can be fully navigated using only a keyboard}{A modal dialog that creates a "Keyboard Trap" from which keyboard-only users cannot escape}{Providing visible focus outlines on active input fields}{Allowing users to pause auto-rotating carousels}

\mcqitem{18}{What is a "Breadcrumb" in digital software navigation?}{A cookie saved on a hard drive}{A secondary navigation trail showing the user's current location within a website hierarchy (e.g., Home > Electronics > Audio)}{A temporary error message}{An advertisement popup}

\mcqitem{19}{What does the Usability heuristic "Visibility of System Status" mean?}{The software must display its server IP address}{The system should always keep users informed about what is going on through appropriate, timely feedback (e.g., loading bars)}{Making the entire source code open source}{Showing employee names on the homepage}

\mcqitem{20}{When an error occurs on a form, what is the best UX practice for displaying error messages?}{Show a cryptic error code like "Error 0x800F" at the top of the page}{Display a clear, friendly text message placed directly adjacent to the offending input field explaining how to fix it}{Clear the entire form and force the user to re-enter all fields}{Hide the submit button permanently}

\mcqitem{21}{Why should destructive actions (like "Delete Account" or "Format Disk") always require explicit confirmation?}{To slow down the user intentionally}{To prevent catastrophic irreversible errors caused by accidental clicks}{To increase database server load}{Because computer code cannot delete data directly}

\mcqitem{22}{What metric measures the percentage of users who successfully finish a designated test task without moderator help?}{Task Completion Rate (Success Rate)}{Net Promoter Score}{Bounce Rate}{Return on Investment}

\mcqitem{23}{Which metric evaluates customer loyalty by asking *"How likely are you to recommend this product to a friend or colleague?"* on a 0–10 scale?}{Customer Effort Score (CES)}{Net Promoter Score (NPS)}{System Usability Scale (SUS)}{Task Completion Time}

\mcqitem{24}{What is a "Single Ease Question" (SEQ) in usability testing?}{A 50-question survey given before the test}{A single 7-point rating scale question asked immediately after a task: *"Overall, how easy or difficult was this task?"*}{A riddle asked by the interviewer}{A multiple-choice exam question}

\mcqitem{25}{What is the main limitation of automated accessibility testing tools (like Lighthouse or Axe)?}{They are too expensive to run}{They can only detect ~30–40\% of accessibility issues automatically (cannot judge text clarity, alt-text quality, or logical flow)}{They destroy web page code}{They only work on mobile phones}

\mcqitem{26}{If a customer loves an airline's mobile booking app but has their luggage lost by baggage handlers, how is the overall Customer Experience (CX) affected?}{CX remains 100\% positive because the app had good usability}{CX is severely degraded because CX reflects the holistic end-to-end journey across digital and physical touchpoints}{CX only measures mobile software interactions}{The airline's Net Promoter Score will automatically rise}

\mcqitem{27}{In digital navigation, what does "Recognition rather than Recall" mean?}{Forcing users to memorize menu paths from previous pages}{Making options, actions, and objects visible so users do not have to remember information across screens}{Using facial recognition to log in}{Asking users to type exact command-line codes}

\mcqitem{28}{What type of touchpoint is a monthly paper electricity bill delivered to a customer's mailbox?}{Digital touchpoint}{Physical touchpoint}{Human touchpoint}{Social touchpoint}

\mcqitem{29}{Why should user testing include participants with diverse abilities and accessibility requirements?}{It is only required if the government threatens legal action}{Inclusive design uncovers barriers early and creates products that are easier and better for everyone}{Accessible testing reduces the need for software code}{Disabled users only test hardware devices}

\mcqitem{30}{What should a design team do immediately after completing a round of user tests?}{Celebrate and launch the product to market without changes}{Synthesize observations, populate a Feedback Capture Matrix, identify critical themes, and iterate the prototype}{Delete the user recordings}{Fire the developers whose features failed}


\newpage
\section*{Answer Key \& Step-by-Step Explanations}
\addcontentsline{toc}{section}{Answer Key \& Explanations}

\begin{answerkeytable}{Unit 5: Quick Answer Key}
\begin{tabular}{c c c c c c c c c c}
\toprule
\textbf{Q1} & \textbf{Q2} & \textbf{Q3} & \textbf{Q4} & \textbf{Q5} & \textbf{Q6} & \textbf{Q7} & \textbf{Q8} & \textbf{Q9} & \textbf{Q10} \\
\textbf{B} & \textbf{A} & \textbf{B} & \textbf{B} & \textbf{B} & \textbf{B} & \textbf{B} & \textbf{B} & \textbf{B} & \textbf{C} \\
\midrule
\textbf{Q11} & \textbf{Q12} & \textbf{Q13} & \textbf{Q14} & \textbf{Q15} & \textbf{Q16} & \textbf{Q17} & \textbf{Q18} & \textbf{Q19} & \textbf{Q20} \\
\textbf{B} & \textbf{B} & \textbf{A} & \textbf{B} & \textbf{C} & \textbf{B} & \textbf{B} & \textbf{B} & \textbf{B} & \textbf{B} \\
\midrule
\textbf{Q21} & \textbf{Q22} & \textbf{Q23} & \textbf{Q24} & \textbf{Q25} & \textbf{Q26} & \textbf{Q27} & \textbf{Q28} & \textbf{Q29} & \textbf{Q30} \\
\textbf{B} & \textbf{A} & \textbf{B} & \textbf{B} & \textbf{B} & \textbf{B} & \textbf{B} & \textbf{B} & \textbf{B} & \textbf{B} \\
\bottomrule
\end{tabular}
\end{answerkeytable}

\begin{expbox}[Detailed Pedagogical Explanations]
\begin{enumerate}[leftmargin=6mm, itemsep=2.5mm]
  \item \textbf{[Q1: Option B]} Testing is an empirical learning tool to discover flaws and validate assumptions with real users.
  \item \textbf{[Q2: Option A]} Usability checks functional interaction quality; validation checks problem-solution alignment.
  \item \textbf{[Q3: Option B]} "Show, Don't Tell" relies on direct user interaction and neutral observation rather than sales pitches.
  \item \textbf{[Q4: Option B]} Neutral probes encourage the user to verbalize their mental model without coaching.
  \item \textbf{[Q5: Option B]} Reassuring participants ensures they don't feel self-conscious or blame themselves for UI failures.
  \item \textbf{[Q6: Option B]} Effective scenarios give realistic goals and context without prescribing button clicks.
  \item \textbf{[Q7: Option B]} Specifying every exact UI step tests obedience, not discoverability or usability.
  \item \textbf{[Q8: Option B]} The think-aloud protocol captures the participant's internal thoughts and expectations live.
  \item \textbf{[Q9: Option B]} The Wishes ($\Delta$) quadrant captures constructive criticisms and unmet needs.
  \item \textbf{[Q10: Option C]} Unanswered questions and user confusion belong in the Questions (?) quadrant.
  \item \textbf{[Q11: Option B]} CX is the overarching journey across all touchpoints; usability is product-specific interaction.
  \item \textbf{[Q12: Option B]} Delivery couriers, support reps, and retail staff are human service touchpoints.
  \item \textbf{[Q13: Option A]} Journey maps systematically track Actions, Thoughts/Emotions, Touchpoints, Pain Points, and Opportunities.
  \item \textbf{[Q14: Option B]} POUR stands for Perceivable, Operable, Understandable, Robust.
  \item \textbf{[Q15: Option C]} WCAG AA mandates at least a 4.5:1 contrast ratio for standard text.
  \item \textbf{[Q16: Option B]} Screen readers read descriptive alternative text (`alt` attribute) aloud for non-visual users.
  \item \textbf{[Q17: Option B]} Keyboard traps trap focus, preventing keyboard-only users from navigating the site.
  \item \textbf{[Q18: Option B]} Breadcrumbs provide spatial orientation and easy navigation back up the hierarchy.
  \item \textbf{[Q19: Option B]} Systems must provide clear and prompt feedback regarding ongoing operations.
  \item \textbf{[Q20: Option B]} Inline error messages adjacent to fields explain the exact issue and recovery step.
  \item \textbf{[Q21: Option B]} Confirmations prevent destructive, irreversible accidents.
  \item \textbf{[Q22: Option A]} Task Completion Rate measures the proportion of successfully achieved goals.
  \item \textbf{[Q23: Option B]} Net Promoter Score measures advocacy and loyalty on an 11-point scale (0–10).
  \item \textbf{[Q24: Option B]} Single Ease Question is a standardized post-task ease rating.
  \item \textbf{[Q25: Option B]} Automated tools catch code syntax but cannot judge human cognitive clarity or task flow.
  \item \textbf{[Q26: Option B]} CX is holistic; a breakdown in physical service destroys overall customer trust.
  \item \textbf{[Q27: Option B]} Minimizing user memory load by making elements visible promotes usability.
  \item \textbf{[Q28: Option B]} Printed bills, packaging, and hardware devices are physical touchpoints.
  \item \textbf{[Q29: Option B]} Inclusive design solves for one and extends to many, creating universal ease of use.
  \item \textbf{[Q30: Option B]} Immediate structured synthesis and iteration turn test observations into product improvements.
\end{enumerate}
\end{expbox}
